\icode{E-A-MU-L2-260626-001}
2021年1月30日，“华龙一号”全球首堆福建福清核电5号机组投入商业运行，这是中国核电发展的重大跨越，标志着我国在三代核电技术领域跻身世界前列，成为继美国、法国、俄罗斯等国家之后真正掌握自主三代核电技术的国家。“华龙一号”运用了先进的百万千瓦级压水堆核电技术，设某核反应方程为\ce{^{235}_{92}U + ^1_0n \to ^{144}_{56}Ba + ^{89}_{36}Kr + 3 ^1_0n}，其中\ce{^{235}_{92}U}、\ce{^{144}_{56}Ba}、\ce{^{89}_{36}Kr}、\ce{^1_0n}的质量分别为\(m_1\)、\(m_2\)、\(m_3\)、\(m_4\)，光速为\(c\)。下列说法正确的是\paren[CD]
\begin{choices}
    \item 此核反应为核聚变反应
    \item \ce{^{144}_{56}Ba}及\ce{^{89}_{36}Kr}的比结合能小于\ce{^{235}_{92}U}的比结合能
    \item 一次核反应释放的核能为\(E =(m_1 - m_2 - m_3 - 2m_4)c^2\)
    \item 核燃料需达到临界体积才能发生链式反应
\end{choices}
\begin{answer}
    CD
\end{answer}